\documentclass[11pt,a4paper]{article}
\usepackage[utf8]{inputenc}
\usepackage[T1]{fontenc}
\usepackage[spanish,es-tabla]{babel}
\usepackage[margin=2cm]{geometry}
\usepackage{graphicx}
\usepackage{tabularx}
\usepackage{longtable}
\usepackage{array}
\usepackage[hidelinks]{hyperref}

\setlength{\parindent}{0pt}
\setlength{\parskip}{6pt}
\renewcommand{\arraystretch}{1.15}
\newcolumntype{Y}{>{\raggedright\arraybackslash}X}

\begin{document}

\begin{titlepage}
\centering

{\bfseries Universidad Técnica Estatal de Quevedo}\\[4pt]
Facultad de Ciencias de la Computación --- Carrera de Ingeniería de Software\\[10pt]

{\bfseries Ingeniería de Requerimientos (ISR-401)}\\
Paralelo B --- SPA 2026--2027\\[10pt]

{\bfseries Evaluación de la Unidad I (EV-PRAC-01)}\\
EV-SUM-01 en el SGA --- Caso «ALERTARÍO»\\[10pt]

Evaluación práctica individual y presencial

\vfill

\noindent\begin{tabular}{@{}ll@{}}
\textbf{Apellidos y nombres:} & ROBINSON ESPINOZA JEANPIERRE\\
\textbf{URL del repositorio:} & \url{https://github.com/jean200525/ISR-401-EV-PRAC-01-}\\
\textbf{Fecha de la sesión:} & 18 de noviembre de 2026\\
\end{tabular}

\vfill

Esta es la transcripción fiel de la prueba escrita a mano el 18 de noviembre de 2026; las fotografías de las hojas están en la carpeta \texttt{evidencias/} del repositorio.

\end{titlepage}

\newpage

\section*{T1}

A.- por lo que vimos en la clase un requisito se lo consideraba una necesidad o condición que debe cumplir un sistema para satisfacer las necesidades de los stakeholders, cada requisito debe estar escrito de manera clara y deben comprobarse.

E7 describe cómo se trabaja actualmente: el vigilante mide el nivel del río a las 6 de la mañana y a las 6 de la tarde, comunica el dato por radio y el operador lo anota en un cuaderno.

E7 no es un requisito del producto ALERTARÍO, porque solamente explica el procedimiento que se utiliza actualmente. Sirve como información para conocer el problema y entender qué actividades se podrían mejorar.

E12 tampoco es un requisito del producto, porque indica cómo se realizará el trabajo durante el proyecto. Es un requisito del proceso y debe registrarse en la planificación del proyecto.

\textbf{T1(b). Diferencia entre StRS, SyRS y SRS}

\noindent\begin{tabularx}{\linewidth}{|c|Y|c|Y|}
\hline
\textbf{Documento} & \textbf{¿Qué describe?} & \textbf{Ejemplo del caso} & \textbf{¿Por qué?}\\
\hline
StRS & Las necesidades y expectativas de los interesados. & E1 & Los presidentes barriales necesitan recibir alertas a tiempo, incluso de madrugada.\\
\hline
SyRS & Los requisitos generales del sistema completo, incluyendo hardware y software. & E3 & Describe cómo las estaciones HL-20 envían las lecturas por LoRaWAN cada 5 minutos.\\
\hline
SRS & Los requisitos específicos que debe cumplir el software. & E5 & Indica que el sistema debe generar informes mensuales en PDF.\\
\hline
\end{tabularx}

\textbf{T1(c). Actividades del proceso de Ingeniería de Requisitos}

\noindent\begin{tabularx}{\linewidth}{|c|c|c|Y|}
\hline
\textbf{ID} & \textbf{Actividad} & \textbf{Tipo} & \textbf{Justificación}\\
\hline
H1 & Elicitación & Central & La analista visita la Unidad de Gestión de Riesgos y conversa con los operadores para conocer sus necesidades.\\
\hline
H2 & Análisis y negociación & Central & Se comparan las alternativas de comunicación y se acuerda utilizar SMS.\\
\hline
H3 & Especificación & Central & La analista escribe los requisitos en un documento y les asigna identificadores.\\
\hline
H4 & Validación & Central & Los interesados revisan el documento y detectan que falta explicar cómo cancelar una alerta.\\
\hline
H5 & Gestión & Transversal & Se registra el cambio de operadora, se revisan los requisitos afectados y se publica una nueva versión.\\
\hline
\end{tabularx}

SWEBOK v4.0 y marco de Ingeniería de Requisitos de Pohl.

\section*{T2. Clasificación y auditoría}

\textbf{T2(a). Clasificación de los enunciados E1 a E12}

\noindent\begin{tabularx}{\linewidth}{|c|Y|Y|}
\hline
\textbf{ID} & \textbf{Tipo} & \textbf{Documento}\\
\hline
E1 & Funcional & StRS\\
\hline
E2 & Requisito de dominio & SyRS\\
\hline
E3 & Restricción técnica & SyRS\\
\hline
E4 & Restricción del proyecto & Ninguno; plan del proyecto\\
\hline
E5 & Funcional & SRS\\
\hline
E6 & No funcional & SyRS\\
\hline
E7 & No es requisito & Ninguno; descripción del proceso actual\\
\hline
E8 & No funcional & SRS\\
\hline
E9 & Restricción técnica & SyRS\\
\hline
E10 & No funcional & SRS\\
\hline
E11 & No funcional & SRS\\
\hline
E12 & Requisito de proceso & Ninguno; plan de requisitos del proyecto\\
\hline
\end{tabularx}

ISO/IEC 29148:2018 y modelo de calidad ISO/IEC 25010:2023.

\textbf{T2(b). Auditoría de los requisitos B1 a B6}

NO LA REALICE

\section*{T3. Redacción de requisitos bien formados}

\noindent\begin{tabularx}{\linewidth}{|c|Y|c|}
\hline
\textbf{ID} & \textbf{Enunciado del requisito} & \textbf{Tipo}\\
\hline
B1a & Cuando existan 500 usuarios conectados al mismo tiempo, el sistema deberá cargar el mapa de las seis estaciones en menos de 3 segundos. & No funcional\\
\hline
B1b & Cuando un vecino se registre por primera vez, la aplicación deberá completar el registro en un máximo de 3 minutos. & No funcional\\
\hline
B2 & Cuando la lectura del nivel del río en el puente principal supere los 4,5 metros, el sistema deberá clasificar la situación como alerta naranja; cuando supere los 5,5 metros, deberá clasificarla como alerta roja. & Funcional / dominio\\
\hline
B3 & Cada vez que las estaciones HL-20 transmitan una lectura, el sistema deberá recibir y registrar el nivel del agua de las seis estaciones cada 5 minutos mediante el protocolo LoRaWAN. & Funcional\\
\hline
B5 & Al finalizar cada mes, el sistema deberá generar un informe en formato PDF con los niveles registrados por cada una de las seis estaciones. & Funcional\\
\hline
N1 & Cuando un operador autorizado cancele una alerta, el sistema deberá registrar el usuario responsable y la acción realizada. & Funcional\\
\hline
\end{tabularx}

\section*{T4. Sustento normativo de la calidad}

\textbf{T4(a). Características de los requisitos individuales según ISO 29148}

\noindent\begin{tabularx}{\linewidth}{|c|Y|c|}
\hline
\textbf{Característica} & \textbf{¿Qué exige?} & \textbf{Propiedad SMART relacionada}\\
\hline
Necesario & Debe representar una necesidad real del sistema o de sus interesados. & Ninguna\\
\hline
Singular & Debe expresar una sola obligación y no mezclar varias condiciones independientes. & Ninguna\\
\hline
No ambiguo & Debe tener un solo significado y utilizar palabras claras. & específico\\
\hline
Verificable & Debe poder comprobarse mediante pruebas, inspecciones o demostraciones. & medible\\
\hline
Factible & Debe poder realizarse con los recursos y condiciones disponibles. & alcanzable y realizable\\
\hline
\end{tabularx}

\textbf{T4(b). Dos defectos del conjunto de requisitos}

\textbf{Defecto 1. Falta de completitud}

\begin{itemize}
\item Característica afectada: Completitud del conjunto de requisitos
\item Identificadores: E10 y H4.
\item Problema: E10 explica que los operadores autorizados pueden emitir o cancelar alertas, pero no detalla cómo debe registrarse correctamente una cancelación. H4 confirma que los interesados encontraron esta falta durante la revisión del documento.
\item Solución: Agregar un requisito que indique que, cuando un operador autorizado cancele una alerta, el sistema registre quién realizó la acción y la información correspondiente.
\end{itemize}

\textbf{Defecto 2. Inconsistencia en las condiciones de funcionamiento}

\begin{itemize}
\item Característica afectada: Consistencia del conjunto de requisitos
\item Identificadores: B3, E3 y E6.
\item Problema: B3 exige que el sistema mida en todos los puntos del río, en todo momento, incluso sin energía ni señal y para siempre. Sin embargo, E3 establece que existen seis estaciones HL-20 que transmiten lecturas por radio cada 5 minutos, mientras que E6 solamente establece una condición de funcionamiento durante cortes eléctricos de hasta 12 horas.
\item Solución: Reemplazar B3 por un requisito realista que establezca la recepción y el registro de las lecturas de las seis estaciones cada 5 minutos, considerando las condiciones de energía y comunicación disponibles.
\end{itemize}

ISO/IEC/IEEE 29148:2018, cláusula 5.2.6.

\end{document}
